%======================================================================
%  Course overview -- a single page, neither chapter nor section
%======================================================================
\cleardoublepage
\phantomsection
\addcontentsline{toc}{chapter}{Course overview}
\markboth{Course overview}{Course overview}
\thispagestyle{plain}

\noindent{\color{gray!70!black}\rule{\textwidth}{0.6pt}}\par\vspace{6pt}\noindent
{\sffamily\bfseries\Huge Course overview\par}
\vspace{18pt}

\paragraph{Nonlinear problems in solid mechanics.}
Nonlinearity enters a mechanical problem in two ways. It may come from the
\emph{material or the structure}: large strains, nonlinear material behaviour
(plasticity, viscoelasticity, viscoplasticity), instability and buckling.
It may also come from the \emph{nature of the question}: optimization and
inverse problems are nonlinear even when the underlying mechanics is linear.
In both cases the common thread is variational: an energy, its derivatives,
and the structure of the equations they produce.

\paragraph{Roadmap.}
\begin{enumerate}[itemsep=2pt]
  \item \emph{Variational formulation.} Energy, weak and strong forms,
  existence, duality, constraints and Lagrange multipliers; the elliptic
  model problems of diffusion and linear elasticity; finite elements.
  \item \emph{Nonlinear material behaviour.} Plasticity, viscoelasticity and
  viscoplasticity: numerical integration schemes (explicit, implicit, radial
  return), their underlying principles, and their link to the thermodynamics
  of generalized standard materials (Kondo, Mandel); large-strain elasticity.
  \item \emph{Stability, optimization and identification.} Free vibrations
  and buckling, plastic instabilities; minimization of cost functionals,
  gradient computations, shape optimization, inverse problems.
\end{enumerate}
The present draft covers the first part, together with a first inverse
problem (crack identification) which already uses the whole variational
toolkit.

\paragraph{What you should be able to do at the end.}
\begin{itemize}[itemsep=2pt]
  \item Understand variational calculus and the correspondence
  $\text{PDE}\leftrightarrow\text{energy}$; formulate the numerics; know where
  to look for the relevant mathematics.
  \item Understand the principles of plasticity, stability and optimization,
  and how a numerical integration scheme actually works.
  \item Model nonlinear material behaviour with computational tools,
  essentially finite elements, with a precise understanding of the concepts
  and good practice.
\end{itemize}

\paragraph{Practical matters.}
Numerics are discussed in Python; the finite element codes used are
FEniCS, Cast3M and Abaqus. No programming is required in advance: running the
provided examples and writing simple loops in Python is enough to start.
Computations should be followed on paper as they are derived, the numerical
examples explored, and the exercises solved. Assessment is by a written exam
(no documents, no computer: short questions and short computations) and a
project with a report and a presentation.

\vspace{10pt}
% ---------------------------------------------------------------------
% How to use this book -- to be written once the layout is settled.
% ---------------------------------------------------------------------
\draftnote{Remark ``How to use this book'' to be added here.}

\cleardoublepage
